\section{The Loop Comes Home}

The backward walk begins with a refusal. The grammar refuses to treat the
return as a memory of the forward path. A memory can be carried on an open
line. It can say that a reader once selected, that a gauge once partitioned,
that a boundary once traced a residue, that a sign once appeared. But a
memory is not yet a loop. To come home, the carried tower must return to the
standing place that first made identity admissible.

The forward tower climbed by earning each permission. Difference came before
sameness. Counting split into whole and part. Finite approach carried
residue. A clock made trial repeatable. A horizon bounded the trial. Truth
collapsed when the bare interior tried to separate affirmation from denial.
The needle restored one bounded identity. The finite gauge read variation.
The invariant became concrete. Charge appeared as oriented residue. The
boundary traced what the interior could not mix. The outside reader oriented
the sign. Relative selection quotiented presentations into a frame. The
native gauge counted what pooled and what selected.

The backward walk asks each of these permissions to return what it borrowed.
The reader returns selection to quotient. The quotient returns orientation to
the boundary that licensed it. The boundary returns charge to loop closure.
The loop returns sign to the invariant that carried it. The invariant returns
activity to the finite gauge. The finite gauge returns stationarity to the
needle, because stationarity needs a place where the record can say that this
reading is the same reading for this purpose. The walk is backward not
because time is reversed, but because dependency is being repaid.

An open descent would not be enough. Suppose the grammar could descend from
reader to boundary, from boundary to invariant, from invariant to gauge, and
then stop just before the needle. It would have a history of how identity
was used, but it would not have restored the right to use identity. The first
distinction would remain a remembered origin rather than a closed support.
The reader could still ask questions, but the questions would float above
the identity that makes an answer repeatable.

Closure forbids that float. A loop comes home only when the final return
identifies with the first sanctioned standing place. The grammar does not
say that the first distinction was always identical with itself in a hidden
storehouse. It says something narrower. After the collapse, one local
identity was manufactured as the needle. The backward walk may close only by
returning to that needle under the same discipline: identity by sanctioned
quotient, not identity by assumption.

What returns is therefore not a copy of the seed. A copy would still ask for
a prior identity relation to compare original and duplicate. The grammar has
not earned that luxury as a primitive. What returns is the same truth in the
only sense available here: the closed reading is quotiented to the needle,
and no permitted test at this level separates the returned mark from the
standing place that first restored truth. The loop reads sameness because the
needle permits it.

This also explains why the untaken arm must remain visible. When truth first
collapsed, affirmation and denial could not be separated from inside the bare
foundation. The needle restored a bounded separator, but it did not erase the
history of the collapse. In the backward walk, the path that would turn
denial into affirmation is carried as a possible variation but is not taken.
The grammar keeps the arm readable so that closure does not become blind.
The loop closes by a lawful return, not by forgetting that another turn was
formally nameable.

The return through the gauge has the same severity. Gates that pooled do not
quietly become selecting gates on the way home. A selected sign does not
become native unless the baseline that selected it is returned and
identified. A boundary trace does not become an interior mixed scalar because
the reader has walked backward past it. Each obligation must be returned in
the form in which it was earned. A borrowed quotient must come back as a
quotient. A carried residue must come back as residue. A loop sign must come
back through the loop.

This is the first closing of the chapter. The tower reversed is not an
erasure of the tower. It is the tower paying its debts in the only order that
can make the final station accountable. When the walk reaches the needle,
the grammar reads the first fact again, but now as the home point of a loop.
The first distinguishable mark has become a return point. The reader at the
end and the needle at the beginning are not two worlds. They are the two
faces of one closed obligation.

Charge can now be counted. An open path could carry a sign or a phase, but
only the closed return has made traversal billable. Once the loop comes home,
the grammar has something to tally: the number of times the carried
distinction has traversed the loop and returned to the needle.
